\documentclass[10pt,a4paper]{amsart}
\usepackage[margin=19mm]{geometry}
\usepackage{amsmath,amssymb,mathtools}
\usepackage[scr=rsfs,cal=euler]{mathalfa}
\usepackage{xfrac}
\usepackage[unicode,hidelinks,bookmarks=false]{hyperref}
\newcommand{\ie}{i.e.\ }
\newcommand{\cf}{cf.\ }
\newcommand{\resp}{respectively\ }
\newcommand{\Subsection}{\subsection}
\providecommand{\nobreakdash}{\nobreak\hbox{-}}
\setlength{\emergencystretch}{2em}
\multlinegap0pt
\usepackage{bm}
\usepackage{bbm}
\usepackage{dsfont}
\usepackage{xspace}
\usepackage{cancel}
\usepackage{mathtools}
\usepackage{amsthm}
\makeatletter
\@ifundefined{theorem}{\newtheorem{theorem}{Theorem}}{}
\@ifundefined{proposition}{\newtheorem{proposition}[theorem]{Proposition}}{}
\@ifundefined{rem}{\newtheorem{rem}[theorem]{Remark}}{}
\makeatother
\DeclarePairedDelimiter{\norma}{\lVert}{\rVert} %norma
\DeclarePairedDelimiter{\valass}{\lvert}{\rvert} %valore assoluto
\newcommand{\prodscal}[2]{\langle{#1} , {#2} \rangle}
\newcommand{\rgrande}{\mathbb{R}}
\title[A model of anisotropic branched optimal transport]{A model of anisotropic branched optimal transport}
\author{Martina Bellettini, Andrea Marchese}
\date{Source: arXiv:2604.17123v1 (CC BY 4.0)}
\begin{document}
\maketitle

\noindent\emph{Source and licence.} A model of anisotropic branched optimal transport. Version arXiv:2604.17123v1 (CC BY 4.0), distributed under the Creative Commons Attribution 4.0 International licence, as recorded in the arXiv licence metadata for this exact version and shown on the abstract page https://arxiv.org/abs/2604.17123v1. Full rights notice: licenses/arxiv-2604.17123-CC-BY-4.0.txt.\par\smallskip

\noindent\textit{Editorial scope.} Counted unit: the Proposition whose source label is \texttt{igperpoligoni} in the article (source0.tex lines 479--492, source label \texttt{igperpoligoni}), delivered together with the article's own complete original proof of it (source0.tex lines 494--602, one \texttt{proof} environment of 109 lines). No mathematics is written, repaired or re-derived: statement and proof are the article's own text line for line. Required context retained uncounted: remverticiefacce (lines 442--478). Every definition, display and label of those lines is retained unchanged. Omitted from this package and not redistributed: 1--441, 493--493, 603--1426. Nothing inside a retained excerpt is omitted or altered: every retained body is delivered exactly as the article's lines are written, so no omission receipt of any kind accompanies this package. Citations: the retained excerpts carry no citation command at all, so no bibliography entry of the article is transcribed and no reference is redistributed. Reference adaptation: none was needed and none was made; every reference inside the delivered unit resolves inside it, so no unresolved reference or citation is produced. Printed slips kept literal: no printed word, symbol, hypothesis or number of the counted statement or of its proof was altered. Layout and class: the article's own class is \texttt{article} with packages including geometry, graphicx, url, booktabs, tabularx, array, float, fancyhdr, inputenc, babel, amsmath, csquotes, hyperref, amsthm; the wrapper is the lane's pinned \texttt{amsart} editorial wrapper at 10pt on A4, which restates the theorem-like environments the retained text needs and adds the article's own macro definitions that the retained text needs (\texttt{\textbackslash norma}, \texttt{\textbackslash prodscal}, \texttt{\textbackslash rgrande}, \texttt{\textbackslash valass}), with extra wrapper packages (bm, bbm, dsfont, xspace, cancel, mathtools). The delivered numbering is the unit's own. Assets: no retained span contains a figure environment, a graphics-inclusion command, a PDF-page inclusion, a table or a TikZ picture; no image file is copied into this package, the package declares no asset dependency and no image file is redistributed with it. Licence: version arXiv:2604.17123v1 of this article is distributed under the Creative Commons Attribution 4.0 International licence, as recorded in the arXiv licence metadata for this exact version and shown on the abstract page https://arxiv.org/abs/2604.17123v1. A full rights notice is delivered at licenses/arxiv-2604.17123-CC-BY-4.0.txt. Characters: every retained excerpt is pure ASCII, so the delivered engine, XeLaTeX with its default Latin Modern text font, reports no missing character.
\begin{rem}\label{remverticiefacce}
    \rm
    \begin{enumerate}
        \item By central symmetry of $P$, for any $i \in \{ 1, \dots, N \}$
        $$
        v_{i+ N} = - v_i, \quad w_{i+ N} = - w_i.
        $$
        \item With this choice of $w_1, \dots, w_{2N}$ we have
        $$
        w_i = \begin{pmatrix}
        w_{i,x} \\
        w_{i,y}
        \end{pmatrix} = \begin{pmatrix}
        -v_{i,y} \\
        v_{i,x}
        \end{pmatrix}
        \text{.}
        $$
        \item If the edge $F_i$ is not parallel to any axis, it lies in a line of equation $ y = m_i x + q_i$ for some nonzero $m_i, q_i \in \rgrande,$ $i \in \{ 1, \dots, 2N \}.$ We have that by symmetry of $P$
        $$
        m_{i +N} = m_i, \quad q_{i + N} = -q_i
        $$
        for any $i \in \{ 1, \dots, N\}.$
        Moreover, a vector $u = \begin{pmatrix}
            u_x \\
            u_y
        \end{pmatrix} \in \rgrande^2$ belongs to the line containing the edge $F_i$ if and only if
        $$
        u_y = m_iu_x + q_i
        \quad
        \iff
        \quad
        a_i u_x + b_i u_y = 1,
        $$
        where $a_i = -\frac{m_i}{q_i}, b_i = \frac{1}{q_i}.$
    \end{enumerate}
\end{rem}

\begin{proposition}\label{igperpoligoni}
Let $P \subset \rgrande^2$ be a centrally symmetric, convex polygon.
Then there exist positive coefficients $\lambda_1, \dots, \lambda_{N} \in \rgrande$ depending on $P$ such that:
\begin{enumerate}
\item\label{enumigperpoligono} for any vector $u \in \partial P$ it holds
\begin{equation}\label{eqigperpoligonipiani}
\sum_{i = 1}^{N} \lambda_i \valass{\prodscal{u}{w_i}} =1,
\end{equation}
\item if $B_r(0) \subset P$ for some $r \in (0, + \infty),$ then there exists a constant $M = M(r)$ such that
    $$
    \sum_{i=1}^{N} \lambda_i \norma{v_i} \leq M.
    $$
\end{enumerate}
\end{proposition}

\begin{proof}
Up to a rotation, we can assume that
\begin{itemize}
    \item none of the vertices of $P$ lies on the coordinate axes,
    \item none of the edges of $P$ is parallel to the coordinate axes.
\end{itemize}
For each $i=1,\dots,2N$, let
$n_i \vcentcolon = (a_i,b_i),$
so that the line containing $F_i$ is given by
$\prodscal{n_i}{u}=1.
$
By Remark \ref{remverticiefacce}, for $i=1,\dots,N$ we have
$n_{i+N}=-n_i$.

We first prove point \ref{enumigperpoligono}. For $k=1,\dots,N$, define $\lambda_k>0$ by
\begin{equation}\label{eq:deflambda_clean}
2\lambda_k w_k = n_k-n_{k-1},
\qquad\text{with }n_0\vcentcolon=-n_N.
\end{equation}
This is well defined, because both $n_k$ and $n_{k-1}$ satisfy
$$
\prodscal{n_k}{v_k}=\prodscal{n_{k-1}}{v_k}=1,
$$
hence
$$
\prodscal{n_k-n_{k-1}}{v_k}=0.
$$
Therefore $n_k-n_{k-1}$ is orthogonal to $v_k$, so it is a scalar multiple of
$w_k=Rv_k$.
Since the outward normals rotate monotonically along the boundary of a convex polygon, this scalar is positive; hence $\lambda_k>0$.

Summing \eqref{eq:deflambda_clean}, for every $k=1,\dots,N$ we obtain
\begin{equation}\label{eq:telescope_clean}
\sum_{i=1}^{k}\lambda_iw_i-\sum_{i=k+1}^{N}\lambda_iw_i=n_k.
\end{equation}
Indeed,
$$
2\sum_{i=1}^{k}\lambda_iw_i=\sum_{i=1}^{k}(n_i-n_{i-1})=n_k+n_N,
$$
while
$$
2\sum_{i=k+1}^{N}\lambda_iw_i=\sum_{i=k+1}^{N}(n_i-n_{i-1})=n_N-n_k,
$$
and subtracting gives \eqref{eq:telescope_clean}.

Now let $u\in F_k$ with $k\in\{1,\dots,N\}$. Since the vertices are ordered counterclockwise and $u$ lies on the edge joining $v_k$ to $v_{k+1}$, we have
$$
\prodscal{u}{w_i}\ge 0 \quad \text{for } i\le k,
\qquad
\prodscal{u}{w_i}\le 0 \quad \text{for } i\ge k+1.
$$
Hence, by \eqref{eq:telescope_clean},
$$
\sum_{i=1}^{N}\lambda_i\valass{\prodscal{u}{w_i}}
=
\prodscal{u}{\sum_{i=1}^{k}\lambda_iw_i-\sum_{i=k+1}^{N}\lambda_iw_i}
=
\prodscal{u}{n_k}
=1.
$$
If instead $u\in F_{k}$ with $k\in\{N+1,\dots,2N\}$, then $-u\in F_{k-N}$ and therefore
$$
\sum_{i=1}^{N}\lambda_i\valass{\prodscal{u}{w_i}}
=
\sum_{i=1}^{N}\lambda_i\valass{\prodscal{-u}{w_i}}
=1.
$$
This proves \eqref{eqigperpoligonipiani}.

We now prove point (2). From \eqref{eq:deflambda_clean} and $\norma{w_i}=\norma{v_i}$ we get
$$
2\lambda_i \norma{v_i} = \norma{n_i-n_{i-1}},
$$
hence
$$
\lambda_i\norma{v_i}
\le
\frac12\bigl(\valass{a_i-a_{i-1}}+\valass{b_i-b_{i-1}}\bigr),
\qquad i=1,\dots,N,
$$
where again $n_0=-n_N$.

Assume now that $B_r(0)\subset P$. Since each line $\prodscal{n_i}{u}=1$ supports $P$, its distance from the origin is at least $r$ and therefore
$$
\norma{n_i}\le \frac1r
\qquad\text{for every }i.
$$
In particular,
$$
\valass{a_i}\le \frac1r,
\qquad
\valass{b_i}\le \frac1r.
$$
For $i=1,\dots,N$, the normals $n_i$ move counterclockwise. After splitting them into at most four sectors, each coordinate is monotone on each sector. Therefore the total variations of the sequences $(a_i)_i$ and $(b_i)_i$ are bounded by
$$
\sum_{i=1}^{N}\valass{a_i-a_{i-1}}\le \frac{8}{r},
\qquad
\sum_{i=1}^{N}\valass{b_i-b_{i-1}}\le \frac{8}{r}.
$$
Summing the previous estimate on $\lambda_i\norma{v_i}$, we conclude
$$
\sum_{i=1}^{N}\lambda_i\norma{v_i}
\le
\frac12
\sum_{i=1}^{N}\bigl(\valass{a_i-a_{i-1}}+\valass{b_i-b_{i-1}}\bigr)
\le \frac{8}{r}.
$$
Hence the thesis holds with $M(r)=8/r$.
\end{proof}

\end{document}
